% fig3_trapping_concept.tex — PDMS vs Sulfur Polymer trap physics (band diagram)
%
% Genre: A (band diagram — mandatory: E-axis, band edges LUMO/HOMO, trap level, 캐리어 위치)
% Reference idiom: band diagram convention from sahaPreventingSpaceCharge2023-02
%   (E-axis vertical + conduction band + valence band + trap levels + n_t/n_m carrier markers)
% Paper verbatim 출처: briefings/concept_briefing.md §3 (LUMO/HOMO/shallow/deep trap 정의),
%   §1 ("심층 트랩을 형성해 코로나 방전으로 주입한 전하를 장기 보유")

\documentclass[border=8pt]{standalone}
\usepackage{polymer-paper-preamble}

\begin{document}
\resizebox{183mm}{!}{%
\begin{tikzpicture}[
  every node/.style={font=\textsf{\fontsize{8}{10}\selectfont}},
  electron/.style={circle, fill=cBlue, draw=cBlue!70!black, inner sep=0pt,
                   minimum size=2.6mm, line width=0.5pt},
  bandLine/.style={line width=1.4pt},
  trapLine/.style={line width=1.1pt, dash pattern=on 2.5pt off 1.5pt},
  injArr/.style={-{Stealth[length=4pt,width=3pt]}, line width=1.0pt},
  axisArr/.style={-{Stealth[length=3.5pt,width=2.5pt]}, cGray, line width=0.6pt},
  recArr/.style={-{Stealth[length=3pt,width=2pt]}, cGray, line width=0.7pt, dashed},
]

% ─── Panel backgrounds ───────────────────────────────────────────────────────
\fill[cLGray!30, rounded corners=2.5mm] (0, 0)    rectangle (5.80, 5.40);
\fill[cLGray!30, rounded corners=2.5mm] (6.20, 0) rectangle (12.00, 5.40);

% ─── Panel labels (a) (b) ───────────────────────────────────────────────────
\node[font=\textsf{\bfseries\fontsize{10}{12}\selectfont}] at (0.30, 5.16) {(a)};
\node[font=\textsf{\bfseries\fontsize{10}{12}\selectfont}] at (6.50, 5.16) {(b)};

% ─── Panel titles ────────────────────────────────────────────────────────────
\node[font=\textsf{\bfseries\fontsize{9.5}{12}\selectfont}] at (2.90, 5.16) {PDMS};
\node[font=\textsf{\fontsize{7.5}{9}\selectfont}, cGray]
    at (2.90, 4.86) {(no trap states in bandgap)};
\node[font=\textsf{\bfseries\fontsize{9.5}{12}\selectfont}] at (9.10, 5.16) {Sulfur Polymer};
\node[font=\textsf{\fontsize{7.5}{9}\selectfont}, cGray]
    at (9.10, 4.86) {(deep trap states)};

% ══════════════════════════════════════════════════════════════════════════════
% PANEL (a): PDMS — clean bandgap (no traps)
% ══════════════════════════════════════════════════════════════════════════════
\begin{scope}[shift={(0.30, 0)}]

  % Conduction band region (light blue tint)
  \fill[cBlue!12] (0.55, 3.80) rectangle (5.10, 4.45);
  % Valence band region (light red tint)
  \fill[cRed!12]  (0.55, 0.75) rectangle (5.10, 1.30);

  % E-axis (vertical) — "E" label at base, anchor=west to right of axis line
  \draw[axisArr] (0.50, 0.55) -- (0.50, 4.55);
  \node[font=\textsf{\fontsize{7}{8.5}\selectfont}, cGray, anchor=south west]
      at (0.55, 4.55) {$E$};

  % LUMO line (top of bandgap, cBlue)
  \draw[bandLine, cBlue] (0.55, 3.80) -- (5.10, 3.80);
  \node[font=\textsf{\fontsize{7}{8.5}\selectfont}, cBlue, anchor=east]
      at (0.42, 3.80) {LUMO};

  % HOMO line (bottom of bandgap, cRed)
  \draw[bandLine, cRed] (0.55, 1.30) -- (5.10, 1.30);
  \node[font=\textsf{\fontsize{7}{8.5}\selectfont}, cRed, anchor=east]
      at (0.42, 1.30) {HOMO};

  % CB / VB labels
  \node[font=\textsf{\fontsize{6.5}{8}\selectfont}, cBlue!70!black]
      at (4.60, 4.18) {CB};
  \node[font=\textsf{\fontsize{6.5}{8}\selectfont}, cRed!70!black]
      at (4.60, 1.00) {VB};

  % Bandgap label (centered in clean gap)
  \node[font=\sffamily\bfseries\fontsize{8}{10}\selectfont, cGray!80!black, align=center]
      at (2.80, 2.80) {clean bandgap\\[1pt]
        {\mdseries\fontsize{6.5}{8}\selectfont(no trap states)}};

  % Mobile electron in CB
  \node[electron] (eA) at (1.30, 4.10) {};
  \draw[injArr, cBlue!70!black] (1.30, 4.10) -- (2.40, 4.10);
  \node[font=\textsf{\fontsize{6.5}{8}\selectfont}, cBlue!80!black, anchor=south]
      at (1.85, 4.20) {mobile};

  % Recombination arrow CB → VB
  \draw[recArr, cGray] (3.80, 3.80) -- (3.80, 1.30);
  \node[font=\textsf{\fontsize{6.5}{8}\selectfont}, cGray!80!black, anchor=west]
      at (3.92, 2.55) {fast};
  \node[font=\textsf{\fontsize{6.5}{8}\selectfont}, cGray!80!black, anchor=west]
      at (3.92, 2.30) {recomb.};

\end{scope}

% ══════════════════════════════════════════════════════════════════════════════
% PANEL (b): Sulfur Polymer — deep trap states in bandgap
% ══════════════════════════════════════════════════════════════════════════════
\begin{scope}[shift={(6.50, 0)}]

  % CB / VB region fills
  \fill[cBlue!12] (0.55, 3.80) rectangle (5.10, 4.45);
  \fill[cRed!12]  (0.55, 0.75) rectangle (5.10, 1.30);

  % E-axis
  \draw[axisArr] (0.50, 0.55) -- (0.50, 4.55);
  \node[font=\textsf{\fontsize{7}{8.5}\selectfont}, cGray, anchor=south west]
      at (0.55, 4.55) {$E$};

  % LUMO + HOMO
  \draw[bandLine, cBlue] (0.55, 3.80) -- (5.10, 3.80);
  \node[font=\textsf{\fontsize{7}{8.5}\selectfont}, cBlue, anchor=east]
      at (0.42, 3.80) {LUMO};
  \draw[bandLine, cRed]  (0.55, 1.30) -- (5.10, 1.30);
  \node[font=\textsf{\fontsize{7}{8.5}\selectfont}, cRed, anchor=east]
      at (0.42, 1.30) {HOMO};

  \node[font=\textsf{\fontsize{6.5}{8}\selectfont}, cBlue!70!black]
      at (4.60, 4.18) {CB};
  \node[font=\textsf{\fontsize{6.5}{8}\selectfont}, cRed!70!black]
      at (4.60, 1.00) {VB};

  % Shallow trap level (just below LUMO, blue dashed)
  \draw[trapLine, cBlue!60!black] (0.85, 3.45) -- (4.50, 3.45);
  \node[font=\textsf{\fontsize{6.5}{8}\selectfont}, cBlue!60!black, anchor=west]
      at (4.55, 3.45) {shallow};

  % Deep trap level (mid-gap, amber, prominent)
  \draw[trapLine, cAmber, line width=1.4pt] (0.85, 2.20) -- (4.50, 2.20);
  \node[font=\textsf{\fontsize{6.5}{8}\selectfont}, cAmber, anchor=west]
      at (4.55, 2.20) {deep ($S_n$)};

  % Trap depth indicator (vertical bracket)
  \draw[cAmber!80!black, line width=0.5pt]
      (0.62, 3.80) -- (0.55, 3.80) -- (0.55, 2.20) -- (0.62, 2.20);
  \node[font=\textsf{\fontsize{6.5}{8}\selectfont}, cAmber!80!black, anchor=east, rotate=90]
      at (0.50, 3.00) {$E_t$};

  % Trapped electrons at deep trap (4 electrons)
  \node[electron] at (1.20, 2.20) {};
  \node[electron] at (2.10, 2.20) {};
  \node[electron] at (2.90, 2.20) {};
  \node[electron] at (3.80, 2.20) {};

  % Trapped electrons at shallow trap (2 electrons)
  \node[electron] at (1.50, 3.45) {};
  \node[electron] at (3.20, 3.45) {};

  % Injected electron entering CB → captured into deep trap
  \node[electron] at (1.20, 4.10) {};
  \draw[injArr, cBlue!70!black] (1.20, 4.05) -- (1.20, 2.30);
  \node[font=\textsf{\fontsize{6.5}{8}\selectfont}, cBlue!80!black, anchor=west]
      at (1.30, 3.10) {capture};

  % Thermal escape (negligible) — small upward arrow with kT << E_t label
  \draw[recArr, cGray!70, line width=0.5pt] (4.10, 2.20) -- (4.10, 3.80);
  \node[font=\textsf{\fontsize{6.5}{8}\selectfont}, cGray!70!black, anchor=west, align=left]
      at (4.20, 3.00) {$kT \ll E_t$\\[-1pt](escape negligible)};

\end{scope}

% ─── Bottom captions (centered within each panel) ────────────────────────────
\node[font=\textsf{\fontsize{7.5}{9}\selectfont}, anchor=north, align=center,
      text width=5.40cm]
    at (2.90, -0.06)
    {empty bandgap $\rightarrow$ injected e$^-$ recombine fast $\rightarrow$ rapid V(t) loss};
\node[font=\textsf{\fontsize{7.5}{9}\selectfont}, anchor=north, align=center,
      text width=5.40cm]
    at (9.10, -0.06)
    {deep traps capture e$^-$ at $E_t$ $\rightarrow$ thermal escape $\ll kT$ $\rightarrow$ long retention};

\end{tikzpicture}%
}
\end{document}
